% RESULTADO_RECUPERADO y exposición autónoma ampliada; no cambia el núcleo común.
% Propietarios cotejados:
% Artículo I REV02/sections/vacancias_capacidad.tex, completo.
% Monografía775/manuscrito/generated/cristal_relojes_sincronizacion.tex,
% líneas1--223 y931--1032: capacidad, retornos e inducción TRIT.
% Las pruebas universales de las dos escalas se reúnen aquí; el comprobador
% ampliacion/verificar_vacancias.py añade un certificado racional finito.
\subsection{Vacancias del refinamiento nonádico y jerarquía de retornos}
\label{vac-add-seccion}

La prolongación APP--TRIT--TPK conserva, junto al bloque expresado, el
cilindro, la orientación y la historia que permiten continuarlo. En su
lectura de capacidad se comparan dos cardinales ya determinados por esa
construcción: las \(3^6=729\) posibilidades de un bloque ternario de seis
posiciones y las \(10^3=1000\) posiciones de su coordenatización decimal de tres
cifras. Esta comparación produce un calendario propio. Algunas
transiciones aumentan el índice de capacidad; otras incorporan un bloque
a la historia manteniendo dicho índice. Estas últimas son las
\emph{vacancias de capacidad} que se estudian a continuación.

El objeto considerado es, por tanto, una representación del refinamiento
conjunto. Su definición conserva el origen en los dos soportes finitos;
la expresión logarítmica permite demostrar las propiedades de su
calendario. El tiempo de bloque, el índice de vacancia y el índice de
capacidad se escribirán por separado. Esta separación hace visible por
qué los retornos \(21/22\) y \(6/7\) pertenecen a escalas sucesivas de
una misma construcción.

\subsubsection{Capacidad decimal y conservación del incremento}
\label{vac-add-capacidad}

\begin{definition}[Índice de capacidad y vacancia]
\label{vac-add-definicion}
En el origen de fase \(0\), se definen
\begin{equation}
 \rho=\log_{1000}729=\log_{10}9,
 \qquad \delta=1-\rho=\log_{10}(10/9),
 \label{vac-add-densidad}
\end{equation}
y, para \(t\in\mathbb Z_{\geq0}\),
\begin{equation}
 \mathcal C_t=\lfloor(t+1)\rho\rfloor,
 \qquad
 z_t=\lfloor(t+1)\delta\rfloor-\lfloor t\delta\rfloor.
 \label{vac-add-indices}
\end{equation}
El entero \(\mathcal C_t\) es el índice de capacidad decimal del lector;
\(z_t=1\) señala una vacancia. La letra \(K\) permanece reservada al
registro dodecafásico y no designa aquí un contador de capacidad.
\end{definition}

Se tiene \(0<\delta<1\). Además, \(\delta\) es irracional: si
\(\rho=a/b\) fuese racional, con \(a,b\) enteros positivos, la igualdad
\(9^b=10^a\) tendría valoración nula en el primo \(5\) a la izquierda
y positiva a la derecha. Esta contradicción demuestra la irracionalidad
de \(\rho\) y de \(1-\rho\). En particular, ningún múltiplo entero
positivo de \(\delta\) pertenece a \(\mathbb Z\), circunstancia que
fija sin ambigüedad los extremos de los intervalos de este desarrollo.

\begin{proposition}[Balance exacto de capacidad y vacancia]
\label{vac-add-balance}
Para \(t\geq1\), \(z_t\in\{0,1\}\) y
\begin{equation}
 \mathcal C_t-\mathcal C_{t-1}=1-z_t.
 \label{vac-add-incremento}
\end{equation}
En consecuencia, para \(0\leq a<b\),
\begin{equation}
 \mathcal C_b-\mathcal C_a+
 \sum_{t=a+1}^{b}z_t=b-a.
 \label{vac-add-conservacion}
\end{equation}
\end{proposition}
\begin{proof}
Como \(0<\delta<1\), la diferencia entre las partes enteras de dos
números separados por \(\delta\) vale \(0\) o \(1\). Por otra parte,
la irracionalidad demostrada da
\[
 \mathcal C_t=t+1-\lceil(t+1)\delta\rceil
             =t-\lfloor(t+1)\delta\rfloor.
\]
La resta de esta igualdad en \(t\) y \(t-1\) prueba
\eqref{vac-add-incremento}; su suma telescópica prueba
\eqref{vac-add-conservacion}.
\end{proof}

Cada transición aporta exactamente una unidad al lado derecho del
balance. Esa unidad se registra como capacidad adquirida cuando
\(z_t=0\), o como vacancia cuando \(z_t=1\). La retención del índice
\(\mathcal C_t\) deja, por ello, una posición recuperable en la
historia. La ecuación expresa conservación del incremento con dos
lecturas, y no una pérdida indeterminada de información.

La densidad también se deduce del mismo balance:
\[
 \frac1N\sum_{t=1}^{N}z_t
 =\frac{\lfloor(N+1)\delta\rfloor}{N}\longrightarrow\delta.
\]
Por tanto, \((z_t)\) no es eventualmente periódica: una sucesión binaria
eventualmente periódica tiene densidad racional. La aperiodicidad
resulta aquí de la regla de capacidad ya fijada por \(729/1000\).

\subsubsection{Determinación exacta de los intervalos \(21/22\)}
\label{vac-add-primarios}

\begin{lemma}[Cotas racionales de la pendiente]
\label{vac-add-cotas}
La densidad de vacancia satisface
\begin{equation}
 \frac7{153}<\delta<\frac6{131}.
 \label{vac-add-cota-delta}
\end{equation}
Si \(\lambda=\delta^{-1}\), \(\beta=22-\lambda\) y
\(\mu=\beta^{-1}\), entonces
\begin{equation}
 21<\lambda<22,
 \qquad \frac17<\beta<\frac16,
 \qquad 6<\mu<7.
 \label{vac-add-cotas-inducidas}
\end{equation}
\end{lemma}
\begin{proof}
Se proporciona una comprobación finita por aritmética racional. Para
\(0<u<1\), sean
\[
 S_N(u)=2\sum_{k=0}^{N}\frac{u^{2k+1}}{2k+1},
 \qquad
 R_N(u)=\frac{2u^{2N+3}}{(2N+3)(1-u^2)}.
\]
La integración de la serie geométrica de \(2/(1-u^2)\) da
\[
 S_N(u)<\log\frac{1+u}{1-u}<S_N(u)+R_N(u).
\]
En efecto, la cola es positiva y, sustituyendo cada denominador
\(2k+1\) por \(2N+3\), queda acotada por la serie geométrica que
define \(R_N\). Se aplican estas desigualdades a
\(u=1/19,1/3,1/9\), que corresponden respectivamente a
\(\log(10/9),\log2,\log(5/4)\), y se utiliza
\(\log10=3\log2+\log(5/4)\).

Con \(N=4\), escribiendo \(A=S_4(1/19)\),
\(B=3S_4(1/3)+S_4(1/9)\) y
\(E=3R_4(1/3)+R_4(1/9)\), se obtiene
\[
 \frac7{153}
 <\frac{A}{B+E}
 <\delta
 <\frac{A+R_4(1/19)}{B}
 <\frac6{131}.
\]
Las dos desigualdades exteriores se comprueban multiplicando por los
denominadores positivos de estas sumas de cinco términos racionales.
Así se prueba \eqref{vac-add-cota-delta} mediante una cota de resto,
sin utilizar una expansión decimal como premisa. Al invertirla se tiene
\(131/6<\lambda<153/7\). Restar a \(22\), y después invertir,
produce todas las desigualdades de \eqref{vac-add-cotas-inducidas}.
\end{proof}

\begin{theorem}[Posiciones y separaciones de las vacancias]
\label{vac-add-teorema-primario}
Las posiciones de los eventos \(z_t=1\), ordenadas desde el primer
evento, son exactamente
\begin{equation}
 p_j=\lfloor j\lambda\rfloor
     =\left\lfloor\frac j\delta\right\rfloor,
 \qquad j\geq1.
 \label{vac-add-posiciones}
\end{equation}
Sus separaciones satisfacen
\begin{equation}
 h_j=p_{j+1}-p_j\in\{21,22\}.
 \label{vac-add-huecos}
\end{equation}
Ambos valores aparecen infinitas veces.
\end{theorem}
\begin{proof}
El evento que atraviesa el entero \(j\) ocurre precisamente cuando
\(t\delta<j<(t+1)\delta\). La irracionalidad excluye igualdad en
los extremos. Dividiendo por \(\delta\) se obtiene
\(t=\lfloor j/\delta\rfloor=p_j\). Como \(\delta<1\), un paso
atraviesa a lo sumo un entero, de modo que la enumeración es exacta.
Escribiendo \(\lambda=21+\eta\), con \(0<\eta<1\), resulta
\[
 h_j=21+\lfloor(j+1)\eta\rfloor-\lfloor j\eta\rfloor,
\]
lo que prueba \eqref{vac-add-huecos}. Finalmente,
\((p_{J+1}-p_1)/J\to\lambda\in(21,22)\). Si uno de los dos valores
apareciese sólo un número finito de veces, esa media convergería al
otro valor entero, contradicción.
\end{proof}

\subsubsection{Inducción sobre los intervalos cortos: los retornos \(6/7\)}
\label{vac-add-secundarios}

\begin{theorem}[Calendario inducido de intervalos cortos]
\label{vac-add-teorema-secundario}
Sea \(s_j=22-h_j\), para \(j\geq1\). Este indicador vale \(1\)
exactamente cuando el intervalo entre vacancias es corto, de longitud
\(21\). Con \(\beta\) y \(\mu\) definidos en el
lema~\ref{vac-add-cotas},
\begin{equation}
 s_j=\lfloor(j+1)\beta\rfloor-\lfloor j\beta\rfloor.
 \label{vac-add-indicador-corto}
\end{equation}
Los índices de los intervalos cortos son
\begin{equation}
 q_m=\lfloor m\mu\rfloor=\lfloor m/\beta\rfloor,
 \qquad q_{m+1}-q_m\in\{6,7\},\quad m\geq1.
 \label{vac-add-posiciones-cortos}
\end{equation}
Las dos separaciones se presentan infinitas veces y su orden no es
eventualmente periódico.
\end{theorem}
\begin{proof}
La irracionalidad de \(\delta\) implica la de \(\lambda\),
\(\beta\) y \(\mu\). Puesto que \(\lambda=22-\beta\), para
\(j\geq1\) se tiene
\[
 p_j=22j-\lceil j\beta\rceil
     =22j-\lfloor j\beta\rfloor-1.
\]
Al restar términos consecutivos aparece
\eqref{vac-add-indicador-corto}. Aplicando al indicador de pendiente
\(\beta\) el argumento de cruce de enteros del teorema anterior,
su evento \(m\) está en \(q_m=\lfloor m/\beta\rfloor\).
La desigualdad \(6<\mu<7\) prueba las dos separaciones posibles.
Su media converge a \(\mu\); por tanto, ambas aparecen infinitas
veces. Además, una sucesión de separaciones eventualmente periódica
tendría media racional, mientras \(\mu\) es irracional. Esto excluye
también una alternancia estricta, cuya media sería \(13/2\).
\end{proof}

Los números \(6\) y \(7\) cuentan aquí \emph{intervalos entre eventos
de vacancia}, no bloques de refinamiento. La vuelta a la escala primaria
es igualmente explícita. Entre \(q_m\) y \(q_{m+1}\) hay un solo
intervalo corto, situado al principio, y todos los demás son largos.
Por consiguiente, si \(d_m=q_{m+1}-q_m\),
\begin{equation}
 p_{q_{m+1}}-p_{q_m}=21+22(d_m-1)=22d_m-1
 \in\{131,153\}.
 \label{vac-add-retorno-bloques}
\end{equation}
Así, \(131=14\cdot9+5\) y \(153=17\cdot9\) describen dos lecturas
de retorno \emph{en el reloj de bloques}. La ecuación conserva la
operación que conecta las escalas, en lugar de identificar sus
cardinales por su sola aparición.

\subsubsection{Capacidad retenida y selección del régimen TRIT}
\label{vac-add-regimen}

Sea \(v_j=\mathcal C_{p_j}\) la capacidad registrada en la vacancia
\(j\). De \(p_j\delta<j<(p_j+1)\delta\) y \(\delta<1\) resulta
\(\lfloor(p_j+1)\delta\rfloor=j\). Por
\eqref{vac-add-indices},
\begin{equation}
 v_j=p_j-j,
 \qquad v_{j+1}-v_j=h_j-1=21-s_j,
 \qquad [v_{j+1}]_3=[v_j-s_j]_3.
 \label{vac-add-transporte-mod3}
\end{equation}
Un intervalo largo conserva la clase módulo \(3\); uno corto la
desplaza una unidad en sentido negativo. Aplicando el selector del
núcleo formal a la fase positiva \(r_j=1+[v_j]_9\), se obtiene
\begin{equation}
 \tau_j=M_{\rm ph}(r_j),\qquad
 M_{\rm ph}(r)=
 \begin{cases}
 +1,&r\in\{1,4,7\},\\
 -1,&r\in\{2,5,8\},\\
 0,&r\in\{3,6,9\}.
 \end{cases}
 \label{vac-add-selector}
\end{equation}
Esta composición fija el régimen inducido por la capacidad en cada
evento. Conserva tanto el orden de los cambios como el número de
eventos durante los cuales permanece cada régimen.

\begin{corollary}[Mesetas del régimen inducido]
\label{vac-add-mesetas}
Para cada \(m\geq1\), la clase \([v_j]_3\), y por tanto
\(\tau_j\), es constante en
\(q_m+1\leq j\leq q_{m+1}\). La longitud de cada una de estas
mesetas completas es \(q_{m+1}-q_m\in\{6,7\}\).
\end{corollary}
\begin{proof}
En el intervalo entero indicado, el siguiente cambio ocurre exactamente
al pasar de \(j=q_{m+1}\) a \(j=q_{m+1}+1\): es la próxima posición
donde \(s_j=1\). En todos los pasos anteriores \(s_j=0\), de modo
que \eqref{vac-add-transporte-mod3} conserva la clase. En el extremo,
la disminuye en una unidad. El selector \eqref{vac-add-selector}
asigna valores distintos a las tres clases; por ello conserva las
mismas mesetas y sus extremos.
\end{proof}

El tramo inicial \(1\leq j\leq q_1\) depende del origen elegido y
es una restricción inicial de una meseta. En el origen presente mide
\(q_1=6\). Las primeras veinte clases son
\[
 \underbrace{2,\ldots,2}_{6},\qquad
 \underbrace{1,\ldots,1}_{7},\qquad
 \underbrace{0,\ldots,0}_{7},
\]
con firmas \(0,-1,+1\), respectivamente. Cada intervalo corto
desencadena el siguiente paso del ciclo de regímenes.

Las dos fases de interés quedan definidas como
\[
 g_t^{\rm cron}=[t]_9,
 \qquad g_t^{\rm cap}=[\mathcal C_t]_9,
 \qquad
 g_t^{\rm cap}-g_{t-1}^{\rm cap}\equiv1-z_t\pmod9.
\]
Una vacancia retiene la fase de capacidad mientras el tiempo de bloque
avanza. La transición del estado enriquecido sigue siendo la composición
TPK de selección, transporte y actualización de memoria. Por ello,
\(\mathcal C_t=\mathcal C_{t-1}\) no reinicializa el estado ni
reemplaza la vuelta holonómica \(\Gamma_9\) por una pausa. También
las dos escalas inducidas mantienen esta distinción: entre los eventos
cortos consecutivos de \eqref{vac-add-retorno-bloques}, el incremento de
capacidad es \(21d_m-1\), es decir, \(125\) o \(146\), mientras
el incremento de bloques es \(131\) o \(153\).

\subsubsection{Muestra racionalmente certificada y procedencia}
\label{vac-add-muestra}

La tabla presenta una muestra del origen canónico. \(h_j\) es la
separación \emph{posterior} a \(p_j\); por ello, la vacancia de
\(j=6\), en el bloque \(131\), inicia el primer intervalo corto.
\begin{center}
\small
\begin{tabular}{r|rrrrrrrrrr}
\toprule
\(j\)&1&2&3&4&5&6&7&8&9&10\\
\midrule
\(p_j\)&21&43&65&87&109&131&152&174&196&218\\
\(h_j\)&22&22&22&22&22&21&22&22&22&22\\
\(v_j\)&20&41&62&83&104&125&145&166&187&208\\
\([v_j]_3\)&2&2&2&2&2&2&1&1&1&1\\
\bottomrule
\end{tabular}
\end{center}
Los primeros índices de intervalos cortos y sus separaciones son
\[
 (q_m)_{m=1}^{10}=(6,13,20,27,34,41,48,54,61,68),
\]
\[
 (q_{m+1}-q_m)_{m=1}^{9}=(7,7,7,7,7,7,6,7,7).
\]
La muestra exhibe el orden efectivo, y las pruebas anteriores fijan
su ley a profundidad arbitraria.

El archivo adjunto \texttt{ampliacion/verificar\_vacancias.py} calcula
horquillas racionales con las series y restos del
lema~\ref{vac-add-cotas}. Cada parte entera se acepta únicamente cuando
los dos extremos de la horquilla producen el mismo entero. Con orden
\(16\) reproduce \(512\) eventos, hasta el bloque \(11189\), y
\(64\) separaciones secundarias, además de las identidades de balance,
los residuos y las mesetas. Esta comprobación finita mantiene separado
su alcance del de los teoremas universales.

El índice de capacidad, su complementariedad con las vacancias y la
inducción del selector TRIT se recuperan del Artículo I y de los
capítulos de calendario y reloj inducido de la monografía sobre orden
topológico-temporal aperiódico. Aquí se reúnen sus definiciones y las
pruebas de ambas escalas con una convención única de índices. El
certificado racional constituye una comprobación adicional de esos
resultados y conserva su procedencia.
